% GENERATED FILE. Edit canonical lessons, then run npm run generate:books.
% canonical-source-hash: fnv1a64:529bab2fc0ad62d7
% canonical-lessons: ML-C385-vivaham, ML-C385-vivahaniscayam, ML-C385-tali, ML-C385-ksanakkattu, ML-C385-vivaham-kalikkuka

\chapter{A Wedding}
\label{ch:ml-a-wedding}
\hlchaptermodality{385}

\begin{chapteropening}
\textbf{By the end of this chapter:} \emph{I can say a marriage, an engagement, a wedding pendant, an invitation card and to marry.}

Complete the last lesson of chapter 385: I can say a marriage, an engagement, a wedding pendant, an invitation card and to marry.
\end{chapteropening}

\section[\texorpdfstring{\ml{വിവാഹം} (vivāhaṁ)}{വിവാഹം (vivāhaṁ)}]{\ml{വിവാഹം} (vivāhaṁ) — a marriage; a wedding}

\label{lesson:ML-C385-vivaham}



\begin{quote}
Before the new one: say the Malayalam for an autobiography, then the Malayalam for a biography.
\end{quote}

\subsection*{You'll want to know: \ml{വിവാഹം}}
\textbf{\ml{വിവാഹം}} — \emph{vivāhaṁ} — \textquotedblleft{}a marriage; a wedding\textquotedblright{}.

\begin{quote}
\ml{അവരുടെ} \ml{വിവാഹം} \ml{അടുത്ത} \ml{മാസമാണ്}.
\end{quote}

\begin{quote}
\emph{avaruṭe vivāhaṁ aṭutta māsamāṇŭ.}
\end{quote}

\begin{quote}
Their wedding is next month.
\end{quote}

The formal word beside \ml{കല്യാണം}; invitations and the news use \ml{വിവാഹം}.

\begin{cousinweb}
\textbf{\ml{വിവാഹം}} (\emph{vivāhaṁ}) is borrowed from Sanskrit \emph{vivāha}, \textquotedblleft{}marriage\textquotedblright{} (CDIAL 11920), from \emph{vi-vah}, \textquotedblleft{}to lead away\textquotedblright{}: the bride was led from her father's house.
\end{cousinweb}

\begin{grammarlens}[title={What you've built}]
One more word for telling what happened.
\end{grammarlens}

\subsection*{Your turn}
\begin{itemize}
\raggedright
  \item \emph{Say it:} \emph{vivāhaṁ}
  \item \emph{Say it:} \emph{vivāhaṁ}, once more
  \item \emph{Say it:} \emph{avaruṭe vivāhaṁ aṭutta māsamāṇŭ.} — Their wedding is next month.
  \item \emph{Recall:} say the Malayalam for an autobiography, then the Malayalam for a biography, then say \emph{vivāhaṁ} again
\end{itemize}

\begin{tcolorbox}[breakable,colback=teal!4,colframe=teal!35!black,title={Before you move on}]
What does \ml{വിവാഹം} mean? (\textquotedblleft{}A marriage; a wedding\textquotedblright{}.) Say it once more.
\end{tcolorbox}
\section[\texorpdfstring{\ml{വിവാഹനിശ്ചയം} (vivāhaniścayaṁ)}{വിവാഹനിശ്ചയം (vivāhaniścayaṁ)}]{\ml{വിവാഹനിശ്ചയം} (vivāhaniścayaṁ) — an engagement}

\label{lesson:ML-C385-vivahaniscayam}



\begin{quote}
Before the new one: say the Malayalam for a biography, then the Malayalam for a marriage.
\end{quote}

\subsection*{You'll want to know: \ml{വിവാഹനിശ്ചയം}}
\textbf{\ml{വിവാഹനിശ്ചയം}} — \emph{vivāhaniścayaṁ} — \textquotedblleft{}an engagement\textquotedblright{}.

\begin{quote}
\ml{ഇന്നലെ} \ml{ചേച്ചിയുടെ} \ml{വിവാഹനിശ്ചയം} \ml{ആയിരുന്നു}.
\end{quote}

\begin{quote}
\emph{innale cēcciyuṭe vivāhaniścayaṁ āyirunnu.}
\end{quote}

\begin{quote}
Yesterday was my sister's engagement.
\end{quote}

Literally the fixing of the marriage: \ml{നിശ്ചയം}, a settling, a decision.

\begin{cousinweb}
\textbf{\ml{വിവാഹനിശ്ചയം}} (\emph{vivāhaniścayaṁ}) is built from two Sanskrit words, \emph{vivāha}, \textquotedblleft{}marriage\textquotedblright{} (CDIAL 11920), and \emph{niścaya}, \textquotedblleft{}ascertainment, fixed opinion\textquotedblright{} (CDIAL 7439).
\end{cousinweb}

\begin{grammarlens}[title={What you've built}]
One more word for telling what happened.
\end{grammarlens}

\subsection*{Your turn}
\begin{itemize}
\raggedright
  \item \emph{Say it:} \emph{vivāhaniścayaṁ}
  \item \emph{Say it:} \emph{vivāhaniścayaṁ}, once more
  \item \emph{Say it:} \emph{innale cēcciyuṭe vivāhaniścayaṁ āyirunnu.} — Yesterday was my sister's engagement.
  \item \emph{Recall:} say the Malayalam for a biography, then the Malayalam for a marriage, then say \emph{vivāhaniścayaṁ} again
\end{itemize}

\begin{tcolorbox}[breakable,colback=teal!4,colframe=teal!35!black,title={Before you move on}]
What does \ml{വിവാഹനിശ്ചയം} mean? (\textquotedblleft{}An engagement\textquotedblright{}.) Say it once more.
\end{tcolorbox}
\section[\texorpdfstring{\ml{താലി} (tāli)}{താലി (tāli)}]{\ml{താലി} (tāli) — a wedding pendant}

\label{lesson:ML-C385-tali}



\begin{quote}
Before the new one: say the Malayalam for a marriage, then the Malayalam for an engagement.
\end{quote}

\subsection*{You'll want to know: \ml{താലി}}
\textbf{\ml{താലി}} — \emph{tāli} — \textquotedblleft{}a wedding pendant\textquotedblright{}.

\begin{quote}
\ml{വിവാഹത്തിന്} \ml{വരൻ} \ml{വധുവിന്റെ} \ml{കഴുത്തിൽ} \ml{താലി} \ml{കെട്ടി}.
\end{quote}

\begin{quote}
\emph{vivāhattinŭ varan vadhuvinṟe kaḻuttil tāli keṭṭi.}
\end{quote}

\begin{quote}
At the wedding the groom tied the tali round the bride's neck.
\end{quote}

Tying it, \ml{താലി} \ml{കെട്ടുക}, is the moment the couple are married in many weddings.

\begin{cousinweb}
\textbf{\ml{താലി}} (\emph{tāli}) is native Dravidian: DEDR 3175 lists Malayalam \emph{tāli}, the centre piece of a neck ornament tied as the badge of marriage, beside Tamil \emph{tāli}.
\end{cousinweb}

\begin{grammarlens}[title={What you've built}]
One more word for telling what happened.
\end{grammarlens}

\subsection*{Your turn}
\begin{itemize}
\raggedright
  \item \emph{Say it:} \emph{tāli}
  \item \emph{Say it:} \emph{tāli}, once more
  \item \emph{Say it:} \emph{vivāhattinŭ varan vadhuvinṟe kaḻuttil tāli keṭṭi.} — At the wedding the groom tied the tali round the bride's neck.
  \item \emph{Recall:} say the Malayalam for a marriage, then the Malayalam for an engagement, then say \emph{tāli} again
\end{itemize}

\begin{tcolorbox}[breakable,colback=teal!4,colframe=teal!35!black,title={Before you move on}]
What does \ml{താലി} mean? (\textquotedblleft{}A wedding pendant\textquotedblright{}.) Say it once more.
\end{tcolorbox}
\section[\texorpdfstring{\ml{ക്ഷണക്കത്ത്} (kṣaṇakkattŭ)}{ക്ഷണക്കത്ത് (kṣaṇakkattŭ)}]{\ml{ക്ഷണക്കത്ത്} (kṣaṇakkattŭ) — an invitation card}

\label{lesson:ML-C385-ksanakkattu}



\begin{quote}
Before the new one: say the Malayalam for an engagement, then the Malayalam for a wedding pendant.
\end{quote}

\subsection*{You'll want to know: \ml{ക്ഷണക്കത്ത്}}
\textbf{\ml{ക്ഷണക്കത്ത്}} — \emph{kṣaṇakkattŭ} — \textquotedblleft{}an invitation card\textquotedblright{}.

\begin{quote}
\ml{വിവാഹത്തിന്റെ} \ml{ക്ഷണക്കത്ത്} \ml{ഇന്ന്} \ml{തപാലിൽ} \ml{വന്നു}.
\end{quote}

\begin{quote}
\emph{vivāhattinṟe kṣaṇakkattŭ innŭ tapālil vannu.}
\end{quote}

\begin{quote}
The wedding invitation came in the post today.
\end{quote}

\ml{ക്ഷണം}, an invitation, + \ml{കത്ത്}, a letter.

\begin{cousinweb}
\textbf{\ml{ക്ഷണക്കത്ത്}} (\emph{kṣaṇakkattŭ}) is a compound of \ml{ക്ഷണം} (\emph{kṣaṇaṁ}), \textquotedblleft{}an invitation\textquotedblright{}, which goes back to Sanskrit \emph{kṣaṇa}, \textquotedblleft{}a moment; a festival\textquotedblright{} (CDIAL 3642), and \ml{കത്ത്} (\emph{kattŭ}), \textquotedblleft{}a letter\textquotedblright{}.
\end{cousinweb}

\begin{grammarlens}[title={What you've built}]
One more word for telling what happened.
\end{grammarlens}

\subsection*{Your turn}
\begin{itemize}
\raggedright
  \item \emph{Say it:} \emph{kṣaṇakkattŭ}
  \item \emph{Say it:} \emph{kṣaṇakkattŭ}, once more
  \item \emph{Say it:} \emph{vivāhattinṟe kṣaṇakkattŭ innŭ tapālil vannu.} — The wedding invitation came in the post today.
  \item \emph{Recall:} say the Malayalam for an engagement, then the Malayalam for a wedding pendant, then say \emph{kṣaṇakkattŭ} again
\end{itemize}

\begin{tcolorbox}[breakable,colback=teal!4,colframe=teal!35!black,title={Before you move on}]
What does \ml{ക്ഷണക്കത്ത്} mean? (\textquotedblleft{}An invitation card\textquotedblright{}.) Say it once more.
\end{tcolorbox}
\section[\texorpdfstring{\ml{വിവാഹം} \ml{കഴിക്കുക} (vivāhaṁ kaḻikkuka)}{വിവാഹം കഴിക്കുക (vivāhaṁ kaḻikkuka)}]{\ml{വിവാഹം} \ml{കഴിക്കുക} (vivāhaṁ kaḻikkuka) — to marry; to get married}

\label{lesson:ML-C385-vivaham-kalikkuka}



\begin{quote}
Before the new one: say the Malayalam for a wedding pendant, then the Malayalam for an invitation card.
\end{quote}

\subsection*{You'll want to know: \ml{വിവാഹം} \ml{കഴിക്കുക}}
\textbf{\ml{വിവാഹം} \ml{കഴിക്കുക}} — \emph{vivāhaṁ kaḻikkuka} — \textquotedblleft{}to marry; to get married\textquotedblright{}.

\begin{quote}
\ml{അവൻ} \ml{ഒരു} \ml{അധ്യാപികയെ} \ml{വിവാഹം} \ml{കഴിച്ചു}.
\end{quote}

\begin{quote}
\emph{avan oru adhyāpikaye vivāhaṁ kaḻiccu.}
\end{quote}

\begin{quote}
He married a teacher.
\end{quote}

Past: \ml{വിവാഹം} \ml{കഴിച്ചു}, \emph{vivāhaṁ kaḻiccu}.

The partner takes -\ml{എ}: \ml{അവളെ} \ml{വിവാഹം} \ml{കഴിച്ചു}. In speech, \ml{കല്യാണം} \ml{കഴിക്കുക}.

\begin{cousinweb}
\textbf{\ml{വിവാഹം} \ml{കഴിക്കുക}} (\emph{vivāhaṁ kaḻikkuka}) joins \ml{വിവാഹം} (\emph{vivāhaṁ}), from Sanskrit \emph{vivāha} (CDIAL 11920), and the native verb \ml{കഴിക്കുക}, \textquotedblleft{}to pass, to get through\textquotedblright{} (DEDR 1356): to get a marriage done.
\end{cousinweb}

\begin{grammarlens}[title={What you've built}]
One more word for telling what happened.
\end{grammarlens}

\subsection*{Your turn}
\begin{itemize}
\raggedright
  \item \emph{Say it:} \emph{vivāhaṁ kaḻikkuka}
  \item \emph{Say it:} \emph{vivāhaṁ kaḻikkuka}, once more
  \item \emph{Say it:} \emph{avan oru adhyāpikaye vivāhaṁ kaḻiccu.} — He married a teacher.
  \item \emph{Recall:} say the Malayalam for a wedding pendant, then the Malayalam for an invitation card, then say \emph{vivāhaṁ kaḻikkuka} again
\end{itemize}

\begin{tcolorbox}[breakable,colback=teal!4,colframe=teal!35!black,title={Before you move on}]
What does \ml{വിവാഹം} \ml{കഴിക്കുക} mean? (\textquotedblleft{}To marry; to get married\textquotedblright{}.) Say it once more.
\end{tcolorbox}
